\chapter{Related Work}
\label{ch:related}

\section{Functional translation: Aeneas}

Aeneas is a verification toolchain for Rust built on the bet that, for
safe code, memory is an implementation detail. Rust sources are compiled
to MIR, lifted by Charon into LLBC (a borrow-explicit intermediate
representation), and then translated into a pure lambda calculus with no
heap, no addresses, and no separation logic~\cite{aeneas-charon}. Shared
references become plain immutable values; mutable references are split
into a forward function (lending the value out) and a backward function
(returning the updated value), so the proof engineer reasons about an
equivalent functional program. Backends target several theorem provers,
including Lean, and the admitted fragment is safe, sequential
Rust~\cite{aeneas-site}.

Circe is Aeneas's closest relative and makes the same bet one level
down: C and C++-lite instead of Rust, ClangIR text instead of MIR/LLBC.
The differences follow from the source language. Rust brings its borrow
checker as a trusted aliasing oracle; C brings no such thing, so Circe
replaces it with an explicit verified gate plus text-derived noalias
facts and a proved memory transfer. Where Aeneas inherits its subset
from the language's safety discipline, Circe defines its subset
document by document and rejects everything else loudly.

\section{MLIR-based analysis: VAST}

VAST, from Trail of Bits, is an experimental compiler pipeline that
lifts C and C++ into a tower of MLIR dialects, one per abstraction
level, so that each analysis can pick the representation that fits
it~\cite{vast-blog}. It is a new compiler middle-end for Clang aimed at
program analysis and instrumentation, not a verifier: it finds bugs, it
does not prove programs correct.

The kinship is at the input layer. Both projects meet C and C++ at MLIR
level rather than at source syntax, and both exploit the structure the
frontend already computed. The goals diverge from there: VAST climbs
down the tower toward analysis-friendly abstractions, while Circe climbs
up from CIR text~\cite{clangir} toward pure functions with proved
equational specs.

\section{Rust verifiers: Prusti and Flux}

Two Rust efforts keep memory in the picture and automate around it.
Prusti is a deductive verifier that embeds annotated Rust into Viper, an
intermediate verification language based on implicit dynamic frames, a
separation-logic variant; the borrow checker's assumptions become
permissions, and verification is automatic~\cite{prusti}. Flux goes the
other way into the type system: it extends Rust with liquid (refinement)
types, tracks strong updates to mutable references by exact location,
and discharges the resulting verification conditions by SMT-backed
inference~\cite{flux}.

Against both, Circe is the more radical erasure. Prusti reasons about
permissions, Flux about refinements; Circe, for admitted shapes, reasons
about neither, because there is no heap left to have permissions over.
The price is the gate: where Prusti and Flux accept broad Rust and ask
for annotations, Circe accepts narrow C/C++ and asks for nothing, since
admission itself is the specification of what may alias what.

\section{Verified compilation: CompCert}

CompCert is the landmark verified C compiler: written and proved in
Rocq (formerly Coq), it comes with a formal semantics of its source
fragments and a proof that compilation preserves it~\cite{compcert}.
Its theorem is about the compiler (the executable behaves like the
source); Circe's theorems are about programs (the emitted function
satisfies its equation). The two compose rather than compete: CompCert
says the toolchain preserves meaning, Circe says what the meaning is
worth. Notably, CompCert's source semantics is a memory semantics ---
precisely the layer Circe erases for the common case.

\section{SMT-backed verification: Dafny}

Dafny is a verification-aware programming language: the programmer
writes code together with specifications, and an SMT-backed pipeline
checks them automatically~\cite{dafny}. It optimizes the edit--specify--
verify loop for new code. Circe optimizes the opposite direction:
existing C and C++ that cannot be rewritten, verified after the fact
through translation rather than through a new language. Dafny's
automation is push-button; Circe's is tactic-backed and interactive past
the simplifier, which is the cost of meeting legacy code where it lives.

\section{Separation-logic frameworks: Iris}

Iris is a higher-order concurrent separation-logic framework embedded in
Rocq, with the MoSeL proof mode as its tactic layer; it is the standard
answer when programs need invariants, ghost state, and framing over
shared mutable state~\cite{iris}. The project evaluated exactly this
option in a time-boxed spike: vendoring an Iris-like framing layer over
the memory model for the one residue class the equational story cannot
cover (conditionally-live writers, where the branch makes footprints
caller-dependent). The verdict was to defer --- the native cost of one
conditional proof stayed small, Iris would not have unblocked admission
by itself, and framing without lifetime facts had nothing to attach to.
Both revisit triggers have since fired --- conditionally-live framing
was admitted natively with proven verdicts, and lifetime facts
arrived --- and the deferral was reaffirmed anyway: no consumer is
waiting on logic-level framing, and native framing keeps winning on
burden. The native path's known price is bespoke posts per conditional
shape, with store equations staying hypotheses and no general frame
rule carrying untouched locations through for free; what a future
vendoring would buy is that reusable rule plus tactic support, once
shapes multiply past the bespoke-post budget. That future has since
arrived in miniature: the vendoring blockers were re-checked and found
gone --- matching toolchain pins, a full core build, working tactics
--- so Iris is now vendored, interface only, with a frame rule proved
once over the memory model and a parity demo recovering the native
conditional conclusions exactly. The standing policy is native first,
Iris on fallback: every problem gets the simple no-Iris attempt first,
and the frame rule waits for the shape whose bespoke post fails ---
the second conditional shape is the named first candidate. Admission
itself is untouched --- no validator, emitter, or transfer churn, and
it replaces neither the oracle nor the memory tags --- so everything
admitted stays natively proved.

\section{The proof assistant: Lean and Mathlib}

Whatever the source language or the translation style, the proofs have
to live somewhere. In Circe they live in Lean, on top of its
mathematical library: the elaborator, kernel, and library tactics are
the automation substrate for every per-shape correctness proof, and
Lean~4 with its pinned Mathlib revision is the only trusted computing
base beyond the compiler front end~\cite{lean4,mathlib4}. That
concentration of trust is deliberate --- it is what lets the oracle
verdicts be a checked cache rather than an axiom, and what makes the
transfer theorems worth stating at all.
